problem:  
Let \( (M^n, g) \) be a **simply connected, compact Riemannian \(n\)-manifold** whose curvature tensor and first covariant derivative satisfy the following strengthened curvature‑homogeneity condition:

> For every pair of points \(p,q\in M\), there exists a linear isometry \(\phi:T_pM\to T_qM\) such that \(\phi^*R_q = R_p\) and \(\phi^*(\nabla R)_q = \nabla R_p\).

In other words, \(M\) is **curvature homogeneous of order one**.  
Assume further that the **sectional curvature satisfies a uniform pinching condition**
\[
0 < \frac{K_{\min}}{K_{\max}} \le \delta < 1
\]
for some constant \(\delta\) independent of the point.

**Question:**  
For which values of \(\delta\) (if any) does a curvature‑homogeneous, but not locally homogeneous, manifold of dimension \(n\ge4\) exist?  
Equivalently, determine whether there exists a critical pinching constant \(\delta_c(n)\) such that if \(\frac{K_{\min}}{K_{\max}} > \delta_c(n)\), then every curvature‑homogeneous manifold of order one must be locally homogeneous.

Why is it a "good" problem:  
This reformulation replaces a possibly redundant concrete case with a general structural conjecture that clarifies an open frontier in Riemannian geometry.  

1. **Sharp and Conceptual Threshold:**  
   The problem asks whether there exists a *quantitative rigidity constant* \(\delta_c(n)\) beyond which curvature‑homogeneity automatically upgrades to local homogeneity. This mirrors classical geometric thresholds like the \(1/4\)-pinching constant for topological rigidity, but in a fundamentally different algebraic setting—curvature jets rather than topology.  

2. **Bridges Separate Theories:**  
   It connects two independent strands of differential geometry:  
   - **Curvature‑homogeneous geometry** (which captures local curvature invariance without full symmetry), and  
   - **Pinching rigidity** (which describes analytic positivity forcing geometric uniformity).  
   Establishing a new “pinching–homogeneity rigidity constant” would relate curvature inequalities to jet‑algebraic invariants for the first time.  

3. **Possible New Field Direction:**  
   If such a \(\delta_c(n)\) exists, it would interpolate between the algebraic threshold of Singer’s theorem and the analytic threshold of the sphere theorem, introducing a quantitative dimension‑dependent invariant linking analysis, geometry, and Lie‑algebraic classification.  

4. **Extreme Simplicity, Deep Consequences:**  
   The question is expressible in one sentence:  
   *“How positive can sectional curvature be before curvature‑homogeneous metrics must be locally homogeneous?”*  
   Yet answering it would require merging the techniques of Ricci‑flow, curvature‑tensor algebra, and classification of homogeneous models—a challenge deep enough for modern research.